% Personal cover letter for the Industrial PhD at SuperFly Quants
% %--------------------------
% % CONFIGURATION
% %--------------------------
\documentclass[11pt,a4paper]{moderncv}
\moderncvstyle{classic}
\moderncvcolor{black}
\definecolor{firstnamecolor}{rgb}{0,0,0}


% Basic packages
\usepackage[utf8]{inputenc}
\usepackage[T1]{fontenc}
\usepackage[hscale=0.7, top=2cm, bottom=2cm]{geometry}
\usepackage[dvipsnames]{xcolor}
\definecolor{PaleBlue}{RGB}{76, 151, 215}
\usepackage[colorlinks=true, urlcolor=PaleBlue, linkcolor=PaleBlue]{hyperref}


\usepackage{microtype}


% Configure font settings
\renewcommand{\rmdefault}{ppl}
\renewcommand{\sfdefault}{phv}
\renewcommand*{\namefont}{\fontsize{26}{18}\mdseries}


% Configure microtype
\microtypesetup{expansion=false}


% Load sensitive information
\input{../../secrets/personal_info_dk.tex}


% %--------------------------
\begin{document}
\makecvtitle
Hello,


\vspace*{2em}
I am applying for the Industrial PhD at SuperFly Quants as the project appealed to me. I like the idea of using structure to make an impossibly large problem tractable. Tensor networks are new to me, but close enough to physics, linear algebra, and machine learning that I can see where my background fits. The combination of a difficult mathematical problem and a real pricing system is particularly appealing.


\hspace*{2em}
I tend to move towards the boundaries between subjects. I studied Engineering Physics, specialised in Statistical Learning, and took a second degree in economics alongside it. Since graduating, I have moved between research, machine learning, software engineering, and project leadership. The path is not perfectly linear, but it has made me comfortable entering unfamiliar fields, asking basic questions, and staying with a problem until I understand both the theory and its implications.


\hspace*{2em}
The work I enjoyed most at RaySearch had this character. I represented complicated patient anatomy in a form that made prediction and optimisation transparent. Though the methods were not the same as used in quantum algorithms, the focus on solving complicated problems, finding useful structure in higher dimensions, then testing whether the model survived contact with reality. Working with medical physicists also taught me that an elegant model is only valuable if its users can understand and trust it.


\hspace*{2em}
Since then, I have spent several years on the less glamorous parts of applied modelling: unclear data, brittle pipelines, domain assumptions, and systems that need to keep working after the initial experiment. I have learned that I enjoy this part too. It is also what attracts me to doing the PhD inside the team building SuperFly rather than as a purely theoretical project.


\hspace*{2em}
Outside work, I climb, spend time in nature, and build things through woodworking and 3D printing. I like hobbies where I can use my creative side as well as my hands. The same combination of curiosity and practical problem-solving is what I enjoy about technical work.


\hspace*{2em}
I am aware that I am not the typical recent graduate described in the listing. I finished my master's in 2020 and have since built a fairly established career. Applying is not an attempt to start over, but I do wish to trade some breadth for depth and spend several years properly understanding one difficult subject. I would enter with more practical experience than most PhD candidates, but also with humility about how much of the theory I still need to learn.


\vspace*{2em}
Thank you for your consideration. I would be happy to discuss the project and whether my somewhat unusual path could be useful to it. Regards,


\vspace*{1em}
Ivar Cashin Eriksson.


\end{document}